\chapter{Short lanes, nested lanes and lane-access capacity}
\label{chap:short-lanes}

\section{Why short lanes require interaction}
A short lane does not provide an independent stop-line saturation-flow stream for the
entire green merely because it has a separate stop line.  Its contribution depends on
physical storage, vehicles already queued in the short lane, continuing supply from the
common upstream approach, signal timing and lane demand.  Short-lane queues can also
affect the adjacent parent lane.

Published short-lane literature describes effects on the short lane, the adjacent
full-length lane or both, depending on queues, timings and lane flows.  PySig therefore
keeps base saturation flow separate from accessible capacity.

\section{Parent hierarchy}
The relationship is defined using \code{parent\_lane\_id}.  Nesting is permitted:
\begin{verbatim}
A1_1
  A1_1L
    A1_1LL
\end{verbatim}
The IDs are conventional only; hierarchy comes from the parent references.

\section{Length definition and storage}
\code{physical\_length\_m} is the actual length of each lane from that lane's start to
the stop line.  Storage is
\begin{equation}
K_i=\frac{L_i}{l_{PCU}}.
\end{equation}
Thus, for a 70 m short lane and 7 m PCU spacing, $K=10$ PCU.  A 35 m nested short lane
has $K=5$ PCU.

Each child must be shorter than its immediate parent.  Parent cycles are prohibited and
parent/child lanes must belong to the same arm and intersection.

\section{Current cycle-average accessible-capacity model}
The current implementation is a transparent cycle-average approximation.  The root
full-length lane supplies sustainable common upstream capacity.  A non-root short lane
can add, per cycle, no more than both its own nominal stop-line capacity and its physical
storage equivalent:
\begin{equation}
\Delta c_i=\min\left(c_{n,i},K_i\frac{3600}{C}\right).
\end{equation}
The total lane-tree capacity is bounded by
\begin{equation}
c_{tree}=\min\left(
\sum_i c_{n,i},
 c_{n,root}+\sum_{i\neq root}\Delta c_i
\right).
\end{equation}

This formulation captures the fact that the finite short lanes can provide additional
initially stored stop-line discharge but cannot act as unlimited independent approach
lanes for the entire cycle.

\section{Allocation across the tree}
The current implementation allocates available tree capacity across lane demand by an
equal-degree-of-saturation water-filling method, subject to each lane's individual
nominal capacity.  The objective is to avoid assigning all of the parent and child
nominal capacities independently when they ultimately rely on common upstream supply.

The method represents a cycle-average capacity interaction.  It does not yet reproduce
second-by-second taper-blockage and access-clearance events.  This distinction should
remain visible in results and methodology metadata.

\section{Nested lanes}
For multiple levels, each node contributes its own storage-limited capacity while the
whole connected tree remains subject to the root/common-upstream constraint.  Solving
the connected tree as one object prevents pairwise parent/child adjustments from double
counting the same upstream constraint.

\section{Monotonic and conservation expectations}
All else equal, increasing usable short-lane length should not reduce total accessible
capacity.  Assigned demand must remain conserved across parent and child lanes, and the
sum of lane accessible capacities must equal the capacity allocated to the tree.

\section{Future event-driven refinement}
A future refinement can use the same hierarchy but explicitly schedule events for
short-lane queue clearance, taper blockage, access restoration and upstream propagation.
This would make the parent/child effect time-dependent while preserving the current YAML
schema and base/accessible-capacity distinction.
